\documentclass[aspectratio=169,12pt]{beamer}

%% ── Packages ──────────────────────────────────────────────────────────────
\usepackage{amsmath}
\usepackage{bm}
\usepackage{mathtools}   % for \underbrace spacing fixes
\usepackage{booktabs}    % for \toprule etc.
\usepackage{xcolor}

%% ── Theme ─────────────────────────────────────────────────────────────────
\usetheme{Madrid}
\usecolortheme{seahorse}
\setbeamertemplate{navigation symbols}{}

%% ── Macros ────────────────────────────────────────────────────────────────
\newcommand{\KL}{\mathrm{KL}}
\newcommand{\tr}{\mathrm{tr}}

%% ══════════════════════════════════════════════════════════════════════════
\begin{document}

\begin{frame}{ELBO Decomposition and the Role of Entropy}

\bigskip

%% ── KL decomposition ──────────────────────────────────────────────────────
The KL divergence between the variational approximation $q$ and the prior $p$
decomposes as:

\bigskip
\begin{equation*}
  \KL(q \| p)
    \;=\;
    \underbrace{-H[q]}_{\displaystyle\text{neg.\ entropy of } q}
    \;-\;
    \underbrace{\mathbb{E}_{q}[\log p(\bm{\beta})]}_{\displaystyle\text{expected log-prior}}
\end{equation*}

\bigskip\bigskip

%% ── ELBO ──────────────────────────────────────────────────────────────────
Maximising the Evidence Lower Bound (ELBO) therefore balances two forces:

\bigskip
\begin{equation*}
  \mathcal{L}(\bm{\theta})
    \;=\;
    \underbrace{\mathbb{E}_{q}[\log p(\mathbf{y} \mid \bm{\beta})]}_{\displaystyle\text{fit the data}}
    \;-\;
    \underbrace{\KL(q \| p)}_{\displaystyle\text{stay close to prior}}
\end{equation*}

\end{frame}

%% ══════════════════════════════════════════════════════════════════════════
\begin{frame}{KL Direction: Collapse vs.\ Expansion}

\textbf{Q:} Why does $\KL(q\|p)$ cause posterior variance \emph{collapse},
while $\KL(p\|q)$ causes posterior variance \emph{expansion}?

\bigskip

The answer is entirely in \textbf{which distribution the expectation is taken under}:

\medskip
\begin{columns}[T]

\column{0.48\textwidth}
\textbf{$\KL(q\|p)$ — used in VI — ``mode-seeking''}
\begin{equation*}
  \mathbb{E}_{\textcolor{blue}{q}}\!\bigl[\log q(\bm{\beta}) - \log p(\bm{\beta}\mid\mathbf{y})\bigr]
\end{equation*}
$q$ samples its \emph{own} regions.
Where $p\approx 0$, $q$ simply avoids that region — no penalty.\\[4pt]
$\Rightarrow$ $q$ \textbf{collapses onto one mode}; tails of $p$ are ignored.

\column{0.48\textwidth}
\textbf{$\KL(p\|q)$ — ``mean-seeking''}
\begin{equation*}
  \mathbb{E}_{\textcolor{red}{p}}\!\bigl[\log p(\bm{\beta}\mid\mathbf{y}) - \log q(\bm{\beta})\bigr]
\end{equation*}
$p$ drives the expectation.
Wherever $p>0$ and $q\approx 0$, $\log q\to-\infty$ — infinite penalty.\\[4pt]
$\Rightarrow$ $q$ must \textbf{cover all of $p$'s mass}; variance expands.

\end{columns}

\bigskip
\begin{center}
\renewcommand{\arraystretch}{1.3}
\begin{tabular}{@{}lll@{}}
\toprule
& \textbf{Penalises} & \textbf{So $q$\ldots} \\
\midrule
$\KL(q\|p)$ & $q>0$ where $p\approx 0$ & shrinks to avoid $p$'s tails $\to$ \textbf{collapse} \\
$\KL(p\|q)$ & $q\approx 0$ where $p>0$ & expands to cover $p$'s tails $\to$ \textbf{expansion} \\
\bottomrule
\end{tabular}
\end{center}

\smallskip
{\small VI uses $\KL(q\|p)$ because sampling from $q$ is tractable;
sampling from the true posterior $p$ is the problem we are trying to solve.}

\end{frame}

\end{document}
